% Shared preamble for DiPECS academic reports.
% Each report should start with:
%   \documentclass[12pt,a4paper]{article}
%   % Shared preamble for DiPECS academic reports.
% Each report should start with:
%   \documentclass[12pt,a4paper]{article}
%   % Shared preamble for DiPECS academic reports.
% Each report should start with:
%   \documentclass[12pt,a4paper]{article}
%   % Shared preamble for DiPECS academic reports.
% Each report should start with:
%   \documentclass[12pt,a4paper]{article}
%   \input{../dipecs-report-preamble}
% and define the metadata macros before \begin{document}.

\usepackage[UTF8]{ctex}
\usepackage{setspace}
\usepackage{amsmath}
\usepackage{graphicx}
\usepackage{booktabs}
\usepackage{array}
\usepackage{tabularx}
\usepackage[a4paper, left=2.5cm, right=2.5cm, top=2.5cm, bottom=2.5cm]{geometry}
\usepackage{microtype}
\usepackage{float}
\usepackage[table]{xcolor}
\usepackage{multirow}
\usepackage{makecell}
\usepackage{caption}
\usepackage{titlesec}
\usepackage{fancyhdr}
\usepackage{lmodern}
\usepackage{lastpage}
\usepackage{enumitem}
\usepackage[export]{adjustbox}
\usepackage{indentfirst}
\usepackage[super,square,sort&compress]{natbib}
\usepackage{hyperref}

\hypersetup{hidelinks}

% Custom column types.
\newcolumntype{P}[1]{>{\raggedright\arraybackslash}p{#1}}
\newcolumntype{Q}[1]{>{\centering\arraybackslash}p{#1}}

\setlength{\headheight}{14pt}
\onehalfspacing
\setstretch{1.5}
\setlength{\parskip}{0.5em}

\pagestyle{fancy}
\fancyhf{}
\fancyhead[L]{\small\kaishu DiPECS 端侧 AIOS 原型研究}
\fancyhead[R]{\small\kaishu \reporttype}
\fancyfoot[C]{\thepage}
\renewcommand{\headrulewidth}{0.4pt}
\renewcommand{\footrulewidth}{0pt}

\titleformat{\section}
  {\color[RGB]{0,0,128}\bfseries\songti\zihao{4}}
  {\thesection}{1em}{}
\titleformat{\subsection}
  {\color[RGB]{0,0,128}\bfseries\songti\zihao{5}}
  {\thesubsection}{1em}{}
\titleformat{\subsubsection}
  {\color[RGB]{0,0,128}\songti\zihao{5}}
  {\thesubsubsection}{1em}{}
\titleformat{\paragraph}
  {\color[RGB]{0,0,128}\bfseries\songti\zihao{-5}}
  {\theparagraph}{1em}{}

% Report metadata (must be redefined in each main file).
\newcommand{\reporttitle}{}
\newcommand{\reportsubtitle}{}
\newcommand{\reporttype}{}
\newcommand{\reportdate}{}

% Cover page macro.
\newcommand{\makecover}{%
  \newgeometry{left=1.75cm, right=1.75cm, top=1.55cm, bottom=1.55cm}%
  \begin{titlepage}
    \thispagestyle{empty}
    \centering
    \vspace*{0.15cm}
    \IfFileExists{../icon.png}{%
      \includegraphics[width=0.40\textwidth]{../icon.png}%
    }{%
      \vspace{2cm}%
    }%

    \vspace{0.7cm}

    {\heiti\zihao{2} \reporttitle \par}
    \ifx\reportsubtitle\empty\else
      \vspace{0.55em}
      {\songti\zihao{4} ——\reportsubtitle \par}
    \fi

    \vspace{0.9cm}
    {\songti\zihao{3} \reporttype \par}

    \vspace{0.85cm}

    {\songti\zihao{4}
    \renewcommand{\arraystretch}{1.5}
    \begin{tabular}{rl}
      项目名称： & \underline{\makebox[7.4cm][c]{DiPECS}} \\
      作者单位： & \underline{\makebox[7.4cm][c]{DiPECS Team}} \\
      报告类型： & \underline{\makebox[7.4cm][c]{\reporttype}} \\
      完成时间： & \underline{\makebox[7.4cm][c]{\reportdate}} \\
    \end{tabular}
    }
  \end{titlepage}
  \restoregeometry
}

% Abstract environment: produces a centered title block; caller supplies
% \noindent\textbf{摘要：} and \noindent\textbf{关键词：} lines.
\newenvironment{dipecsabstract}{%
  \begin{center}
    {\heiti\zihao{2} \reporttitle \par}
    \ifx\reportsubtitle\empty\else
      \vspace{0.45em}
      {\songti\zihao{4} ——\reportsubtitle \par}
    \fi
  \end{center}
  \vspace{1em}
}{%
  \vspace{1em}
}

% Lightweight table rule used between data rows.
\newcommand{\dtr}{\specialrule{0.25pt}{0.25em}{0.25em}}

% Standard DiPECS table environment: [H], centered, small, caption above,
% label, and consistent row stretch.
\newenvironment{dipecstable}[2]{%
  \begin{table}[H]
    \centering
    \small
    \caption{#1}
    \label{#2}
    \renewcommand{\arraystretch}{1.3}
}{%
  \end{table}
}

% and define the metadata macros before \begin{document}.

\usepackage[UTF8]{ctex}
\usepackage{setspace}
\usepackage{amsmath}
\usepackage{graphicx}
\usepackage{booktabs}
\usepackage{array}
\usepackage{tabularx}
\usepackage[a4paper, left=2.5cm, right=2.5cm, top=2.5cm, bottom=2.5cm]{geometry}
\usepackage{microtype}
\usepackage{float}
\usepackage[table]{xcolor}
\usepackage{multirow}
\usepackage{makecell}
\usepackage{caption}
\usepackage{titlesec}
\usepackage{fancyhdr}
\usepackage{lmodern}
\usepackage{lastpage}
\usepackage{enumitem}
\usepackage[export]{adjustbox}
\usepackage{indentfirst}
\usepackage[super,square,sort&compress]{natbib}
\usepackage{hyperref}

\hypersetup{hidelinks}

% Custom column types.
\newcolumntype{P}[1]{>{\raggedright\arraybackslash}p{#1}}
\newcolumntype{Q}[1]{>{\centering\arraybackslash}p{#1}}

\setlength{\headheight}{14pt}
\onehalfspacing
\setstretch{1.5}
\setlength{\parskip}{0.5em}

\pagestyle{fancy}
\fancyhf{}
\fancyhead[L]{\small\kaishu DiPECS 端侧 AIOS 原型研究}
\fancyhead[R]{\small\kaishu \reporttype}
\fancyfoot[C]{\thepage}
\renewcommand{\headrulewidth}{0.4pt}
\renewcommand{\footrulewidth}{0pt}

\titleformat{\section}
  {\color[RGB]{0,0,128}\bfseries\songti\zihao{4}}
  {\thesection}{1em}{}
\titleformat{\subsection}
  {\color[RGB]{0,0,128}\bfseries\songti\zihao{5}}
  {\thesubsection}{1em}{}
\titleformat{\subsubsection}
  {\color[RGB]{0,0,128}\songti\zihao{5}}
  {\thesubsubsection}{1em}{}
\titleformat{\paragraph}
  {\color[RGB]{0,0,128}\bfseries\songti\zihao{-5}}
  {\theparagraph}{1em}{}

% Report metadata (must be redefined in each main file).
\newcommand{\reporttitle}{}
\newcommand{\reportsubtitle}{}
\newcommand{\reporttype}{}
\newcommand{\reportdate}{}

% Cover page macro.
\newcommand{\makecover}{%
  \newgeometry{left=1.75cm, right=1.75cm, top=1.55cm, bottom=1.55cm}%
  \begin{titlepage}
    \thispagestyle{empty}
    \centering
    \vspace*{0.15cm}
    \IfFileExists{../icon.png}{%
      \includegraphics[width=0.40\textwidth]{../icon.png}%
    }{%
      \vspace{2cm}%
    }%

    \vspace{0.7cm}

    {\heiti\zihao{2} \reporttitle \par}
    \ifx\reportsubtitle\empty\else
      \vspace{0.55em}
      {\songti\zihao{4} ——\reportsubtitle \par}
    \fi

    \vspace{0.9cm}
    {\songti\zihao{3} \reporttype \par}

    \vspace{0.85cm}

    {\songti\zihao{4}
    \renewcommand{\arraystretch}{1.5}
    \begin{tabular}{rl}
      项目名称： & \underline{\makebox[7.4cm][c]{DiPECS}} \\
      作者单位： & \underline{\makebox[7.4cm][c]{DiPECS Team}} \\
      报告类型： & \underline{\makebox[7.4cm][c]{\reporttype}} \\
      完成时间： & \underline{\makebox[7.4cm][c]{\reportdate}} \\
    \end{tabular}
    }
  \end{titlepage}
  \restoregeometry
}

% Abstract environment: produces a centered title block; caller supplies
% \noindent\textbf{摘要：} and \noindent\textbf{关键词：} lines.
\newenvironment{dipecsabstract}{%
  \begin{center}
    {\heiti\zihao{2} \reporttitle \par}
    \ifx\reportsubtitle\empty\else
      \vspace{0.45em}
      {\songti\zihao{4} ——\reportsubtitle \par}
    \fi
  \end{center}
  \vspace{1em}
}{%
  \vspace{1em}
}

% Lightweight table rule used between data rows.
\newcommand{\dtr}{\specialrule{0.25pt}{0.25em}{0.25em}}

% Standard DiPECS table environment: [H], centered, small, caption above,
% label, and consistent row stretch.
\newenvironment{dipecstable}[2]{%
  \begin{table}[H]
    \centering
    \small
    \caption{#1}
    \label{#2}
    \renewcommand{\arraystretch}{1.3}
}{%
  \end{table}
}

% and define the metadata macros before \begin{document}.

\usepackage[UTF8]{ctex}
\usepackage{setspace}
\usepackage{amsmath}
\usepackage{graphicx}
\usepackage{booktabs}
\usepackage{array}
\usepackage{tabularx}
\usepackage[a4paper, left=2.5cm, right=2.5cm, top=2.5cm, bottom=2.5cm]{geometry}
\usepackage{microtype}
\usepackage{float}
\usepackage[table]{xcolor}
\usepackage{multirow}
\usepackage{makecell}
\usepackage{caption}
\usepackage{titlesec}
\usepackage{fancyhdr}
\usepackage{lmodern}
\usepackage{lastpage}
\usepackage{enumitem}
\usepackage[export]{adjustbox}
\usepackage{indentfirst}
\usepackage[super,square,sort&compress]{natbib}
\usepackage{hyperref}

\hypersetup{hidelinks}

% Custom column types.
\newcolumntype{P}[1]{>{\raggedright\arraybackslash}p{#1}}
\newcolumntype{Q}[1]{>{\centering\arraybackslash}p{#1}}

\setlength{\headheight}{14pt}
\onehalfspacing
\setstretch{1.5}
\setlength{\parskip}{0.5em}

\pagestyle{fancy}
\fancyhf{}
\fancyhead[L]{\small\kaishu DiPECS 端侧 AIOS 原型研究}
\fancyhead[R]{\small\kaishu \reporttype}
\fancyfoot[C]{\thepage}
\renewcommand{\headrulewidth}{0.4pt}
\renewcommand{\footrulewidth}{0pt}

\titleformat{\section}
  {\color[RGB]{0,0,128}\bfseries\songti\zihao{4}}
  {\thesection}{1em}{}
\titleformat{\subsection}
  {\color[RGB]{0,0,128}\bfseries\songti\zihao{5}}
  {\thesubsection}{1em}{}
\titleformat{\subsubsection}
  {\color[RGB]{0,0,128}\songti\zihao{5}}
  {\thesubsubsection}{1em}{}
\titleformat{\paragraph}
  {\color[RGB]{0,0,128}\bfseries\songti\zihao{-5}}
  {\theparagraph}{1em}{}

% Report metadata (must be redefined in each main file).
\newcommand{\reporttitle}{}
\newcommand{\reportsubtitle}{}
\newcommand{\reporttype}{}
\newcommand{\reportdate}{}

% Cover page macro.
\newcommand{\makecover}{%
  \newgeometry{left=1.75cm, right=1.75cm, top=1.55cm, bottom=1.55cm}%
  \begin{titlepage}
    \thispagestyle{empty}
    \centering
    \vspace*{0.15cm}
    \IfFileExists{../icon.png}{%
      \includegraphics[width=0.40\textwidth]{../icon.png}%
    }{%
      \vspace{2cm}%
    }%

    \vspace{0.7cm}

    {\heiti\zihao{2} \reporttitle \par}
    \ifx\reportsubtitle\empty\else
      \vspace{0.55em}
      {\songti\zihao{4} ——\reportsubtitle \par}
    \fi

    \vspace{0.9cm}
    {\songti\zihao{3} \reporttype \par}

    \vspace{0.85cm}

    {\songti\zihao{4}
    \renewcommand{\arraystretch}{1.5}
    \begin{tabular}{rl}
      项目名称： & \underline{\makebox[7.4cm][c]{DiPECS}} \\
      作者单位： & \underline{\makebox[7.4cm][c]{DiPECS Team}} \\
      报告类型： & \underline{\makebox[7.4cm][c]{\reporttype}} \\
      完成时间： & \underline{\makebox[7.4cm][c]{\reportdate}} \\
    \end{tabular}
    }
  \end{titlepage}
  \restoregeometry
}

% Abstract environment: produces a centered title block; caller supplies
% \noindent\textbf{摘要：} and \noindent\textbf{关键词：} lines.
\newenvironment{dipecsabstract}{%
  \begin{center}
    {\heiti\zihao{2} \reporttitle \par}
    \ifx\reportsubtitle\empty\else
      \vspace{0.45em}
      {\songti\zihao{4} ——\reportsubtitle \par}
    \fi
  \end{center}
  \vspace{1em}
}{%
  \vspace{1em}
}

% Lightweight table rule used between data rows.
\newcommand{\dtr}{\specialrule{0.25pt}{0.25em}{0.25em}}

% Standard DiPECS table environment: [H], centered, small, caption above,
% label, and consistent row stretch.
\newenvironment{dipecstable}[2]{%
  \begin{table}[H]
    \centering
    \small
    \caption{#1}
    \label{#2}
    \renewcommand{\arraystretch}{1.3}
}{%
  \end{table}
}

% and define the metadata macros before \begin{document}.

\usepackage[UTF8]{ctex}
\usepackage{setspace}
\usepackage{amsmath}
\usepackage{graphicx}
\usepackage{booktabs}
\usepackage{array}
\usepackage{tabularx}
\usepackage[a4paper, left=2.5cm, right=2.5cm, top=2.5cm, bottom=2.5cm]{geometry}
\usepackage{microtype}
\usepackage{float}
\usepackage[table]{xcolor}
\usepackage{multirow}
\usepackage{makecell}
\usepackage{caption}
\usepackage{titlesec}
\usepackage{fancyhdr}
\usepackage{lmodern}
\usepackage{lastpage}
\usepackage{enumitem}
\usepackage[export]{adjustbox}
\usepackage{indentfirst}
\usepackage[super,square,sort&compress]{natbib}
\usepackage{hyperref}

\hypersetup{hidelinks}

% Custom column types.
\newcolumntype{P}[1]{>{\raggedright\arraybackslash}p{#1}}
\newcolumntype{Q}[1]{>{\centering\arraybackslash}p{#1}}

\setlength{\headheight}{14pt}
\onehalfspacing
\setstretch{1.5}
\setlength{\parskip}{0.5em}

\pagestyle{fancy}
\fancyhf{}
\fancyhead[L]{\small\kaishu DiPECS 端侧 AIOS 原型研究}
\fancyhead[R]{\small\kaishu \reporttype}
\fancyfoot[C]{\thepage}
\renewcommand{\headrulewidth}{0.4pt}
\renewcommand{\footrulewidth}{0pt}

\titleformat{\section}
  {\color[RGB]{0,0,128}\bfseries\songti\zihao{4}}
  {\thesection}{1em}{}
\titleformat{\subsection}
  {\color[RGB]{0,0,128}\bfseries\songti\zihao{5}}
  {\thesubsection}{1em}{}
\titleformat{\subsubsection}
  {\color[RGB]{0,0,128}\songti\zihao{5}}
  {\thesubsubsection}{1em}{}
\titleformat{\paragraph}
  {\color[RGB]{0,0,128}\bfseries\songti\zihao{-5}}
  {\theparagraph}{1em}{}

% Report metadata (must be redefined in each main file).
\newcommand{\reporttitle}{}
\newcommand{\reportsubtitle}{}
\newcommand{\reporttype}{}
\newcommand{\reportdate}{}

% Cover page macro.
\newcommand{\makecover}{%
  \newgeometry{left=1.75cm, right=1.75cm, top=1.55cm, bottom=1.55cm}%
  \begin{titlepage}
    \thispagestyle{empty}
    \centering
    \vspace*{0.15cm}
    \IfFileExists{../icon.png}{%
      \includegraphics[width=0.40\textwidth]{../icon.png}%
    }{%
      \vspace{2cm}%
    }%

    \vspace{0.7cm}

    {\heiti\zihao{2} \reporttitle \par}
    \ifx\reportsubtitle\empty\else
      \vspace{0.55em}
      {\songti\zihao{4} ——\reportsubtitle \par}
    \fi

    \vspace{0.9cm}
    {\songti\zihao{3} \reporttype \par}

    \vspace{0.85cm}

    {\songti\zihao{4}
    \renewcommand{\arraystretch}{1.5}
    \begin{tabular}{rl}
      项目名称： & \underline{\makebox[7.4cm][c]{DiPECS}} \\
      作者单位： & \underline{\makebox[7.4cm][c]{DiPECS Team}} \\
      报告类型： & \underline{\makebox[7.4cm][c]{\reporttype}} \\
      完成时间： & \underline{\makebox[7.4cm][c]{\reportdate}} \\
    \end{tabular}
    }
  \end{titlepage}
  \restoregeometry
}

% Abstract environment: produces a centered title block; caller supplies
% \noindent\textbf{摘要：} and \noindent\textbf{关键词：} lines.
\newenvironment{dipecsabstract}{%
  \begin{center}
    {\heiti\zihao{2} \reporttitle \par}
    \ifx\reportsubtitle\empty\else
      \vspace{0.45em}
      {\songti\zihao{4} ——\reportsubtitle \par}
    \fi
  \end{center}
  \vspace{1em}
}{%
  \vspace{1em}
}

% Lightweight table rule used between data rows.
\newcommand{\dtr}{\specialrule{0.25pt}{0.25em}{0.25em}}

% Standard DiPECS table environment: [H], centered, small, caption above,
% label, and consistent row stretch.
\newenvironment{dipecstable}[2]{%
  \begin{table}[H]
    \centering
    \small
    \caption{#1}
    \label{#2}
    \renewcommand{\arraystretch}{1.3}
}{%
  \end{table}
}
